% Source: https://artofproblemsolving.com/community/c3h1657695_2018_ms_mte_theme_round_all
% Stable identity: aops-topic-1657695-q14
If $f$ has roots $r_1,\dots,r_5$, then $g(x)=f(x+4)=0$ exactly when $x+4=r_i$, so $g$ has roots $r_i-4$. By Vieta's formulas $\sum r_i=-\dfrac{1}{2}$, so the sum of the roots of $g$ is $-\dfrac{1}{2}-5\cdot4=\boxed{-\dfrac{41}{2}}$. (In fact $f(x)=(2x+1)(x^2-4)(x^2-9)$ with roots $-\dfrac{1}{2},\pm2,\pm3$.)
